\documentclass[11pt]{article}
% Shared preamble of the RMC kernel write-ups (\input{preamble} after \documentclass).
\usepackage[margin=1in]{geometry}
\usepackage{amsmath,amssymb,bm}
\usepackage{graphicx}
\usepackage{booktabs}
\usepackage{hyperref}

\newcommand{\dd}{\,\mathrm{d}}
\newcommand{\Ein}{E}
\newcommand{\Eout}{E'}
\newcommand{\Ecm}{E_{c}}
\newcommand{\mucm}{\mu_{c}}
\newcommand{\mulab}{\mu_0}

\newcommand{\acespec}[1]{J.~L.~Conlin and P.~Romano,
\href{https://mcnp-green.lanl.gov/pdf_files/TechReport_2019_LANL_LA-UR-19-29016_ConlinRomano.pdf}
{\emph{A Compact ENDF (ACE) Format Specification}}, LA-UR-19-29016 (2019),
Sec.~#1}
\newcommand{\ENDFspec}[1]{\href{https://www.nndc.bnl.gov/endfdocs/ENDF-102-2023.pdf}
{\emph{ENDF-6 Formats Manual}} (2023), #1}

% Common closing section: how the verification figures are produced
\newcommand{\verificationmethod}{%
Each figure is produced by \texttt{verification/verify\_laws.py}. For an incident
energy $\Ein$ along $+z$, $10^6$ collisions are sampled with MC/DC's own reaction
sampler (\texttt{sample\_elastic\_scattering}, \texttt{sample\_inelastic\_scattering}
or \texttt{sample\_fission}), and every emitted neutron is histogrammed in
$48 \times 24$ bins of $(\Eout, \mulab)$, with $\mulab$ the lab cosine to $+z$. The
inverted kernel used by RMC is integrated over the same bins (Gauss--Legendre panels;
two-body lines are split where $\mulab^*(\Eout)$ crosses a cosine edge), multiplied by
the number of collisions, and compared with a Pearson-type $\chi^2$ diagnostic
over bins with at least 5 expected counts (the rest pooled into one bin).
The existing tables list the script's $\chi^2/\mathrm{dof}$ and nominal $p$-value,
the emitted neutrons per collision (sampled vs.\ the integral of the
inverted kernel), and the largest relative change of a bin integral when the
quadrature panels are doubled. Left panels: outgoing-energy marginal; middle:
lab-cosine marginal; right: per-bin $z = (\text{sampled} - \text{inverted}) /
\sqrt{\text{inverted}}$.

\noindent These figures check agreement with the MC/DC sampler, not independent
physical validity. Their existing $p$-values are not calibrated goodness-of-fit
probabilities: the script uses $B$ degrees of freedom, whereas a fixed-total
multinomial histogram with $B$ categories and no fitted parameters uses $B-1$.
Random multiplicities and spectrum assignments can also change the collision-level
covariance from independent Poisson counts; $z$ is then only a count diagnostic.
See \href{https://docs.scipy.org/doc/scipy/reference/generated/scipy.stats.chisquare.html}
{the Pearson test's sampling assumptions}.
Energy bins use the sample extrema, so unsampled tails are excluded and the
reported inverted yield is an integral over that observed window. A formal test
needs independently specified bins, explicit tail categories, and a covariance or
resampling model appropriate to independent collisions. Zero expected density
with observed emissions must be reported rather than discarded. Doubling quadrature
panels gives a convergence indicator, not a rigorous error bound. The stored
figures and tables have not been regenerated for these qualifications.}

\usepackage{amsthm}

\newtheorem{proposition}{Proposition}
\newcommand{\tpsi}{\tilde\psi}
\newcommand{\teps}{\tilde\epsilon}
\newcommand{\Sbar}{\bar S}
\newcommand{\Sigt}{\Sigma_t}
\newcommand{\Om}{\boldsymbol\Omega}
\newcommand{\tb}{\mathbf t}
\newcommand{\E}{\mathbb{E}}

\title{Polynomial spatial bases for the 3D Cartesian RMC solver}
\date{}

\begin{document}
\maketitle

\noindent This note gives the mathematics of the linear and quadratic spatial trial
spaces of the 3D solver (\texttt{mcdc/rmc/space3d.py}, \texttt{polynomial3d.py},
\texttt{driver3d.run3d} with \texttt{spatial\_basis}), and of the MC/DC tally scores
they are projected from (\texttt{flux-xyz<a><b><c>}). Notation follows
\texttt{cartesian\_3d/cartesian\_3d.tex} and
\texttt{cartesian\_3d/tabulated\_correction.tex}.

\section{Why}

With $\tpsi$ constant on each cell, the residual of the collision-only fixed point
does not vanish as the noise does. On the C5G7-style pin (thermal window, cells of
0.15--1\,cm), the collision and face parts of $\|r\|_1$ stayed at 0.51 and 0.60 of
$\|q\|_1$ when the histories per iteration grew tenfold (\texttt{NOTES.md}, section
13). Streaming enters a constant-in-space trial function only through the jumps
across faces, and those are of the order of the flux variation across a cell. Each
iteration and correction pass then has about the per-history variance of a
standard Monte Carlo solve. A spatial basis that represents the in-cell variation
reduces this structural residual: the face jumps of a discontinuous basis of degree
$p$ scale like $h^{p+1}$ for a smooth flux, and a continuous basis has none.

\section{Trial spaces}

On cell $V_k$ with local coordinates $\tb=(t_x,t_y,t_z)\in[-1,1]^3$,
$x = x_k + (t_x+1)h_x/2$, and for each energy bin $g$ and patch $j$ (both still
piecewise constant),
\begin{equation}
  \tpsi(\tb) = \sum_{m} C_m\, P_a(t_x)P_b(t_y)P_c(t_z),\qquad m=(a,b,c),
\end{equation}
with $P_n$ the Legendre polynomials [1, Ch.~22]. Every basis is held in this form,
so the residual, its sampling and the correction passes are written once.
\begin{itemize}
  \item \texttt{dg1}, \texttt{dg2}: discontinuous, complete polynomials,
  $a+b+c\le p$ (4 and 10 modes), the discontinuous Galerkin spaces of [2--4].
  \item \texttt{cg1}, \texttt{cg2}: continuous tensor-product Lagrange spaces $Q_p$
  [5]. Nodes sit at the cell vertices, and for $p=2$ also at the edge, face and
  cell midpoints. The unknowns are nodal values. On a cell, the 1D Lagrange functions
  in Legendre form are
  \begin{equation}
    p=1:\ \phi_{0,1} = \tfrac12 P_0 \mp \tfrac12 P_1;\qquad
    p=2:\ \phi_{0,2} = \tfrac16 P_0 \mp \tfrac12 P_1 + \tfrac13 P_2,\quad
    \phi_1 = \tfrac23 P_0 - \tfrac23 P_2,
  \end{equation}
  and the tensor products give the cell coefficients $C_m$ ($a,b,c\le p$) of the
  nodal values exactly.
\end{itemize}
The constant basis is the case $p=0$ and keeps its own code path.

\section{Projection from tally moments}

\paragraph{Moments.} The update is the $L^2$ projection of $\teps$. On each cell,
energy bin and patch it needs the moments $\int_{V_k}\teps\,P_m$. The new MC/DC
track-length scores \texttt{flux-xyz<a><b><c>} estimate them: a track segment
$\mathbf r_0+s\Om$, $s\in[0,d]$, inside a cell contributes $w\int_0^d
P_aP_bP_c\,\mathrm ds$. This is the track-length estimator of a functional expansion
tally [6]. The integrand is a polynomial of degree $a+b+c\le 6$ in $s$, so a
4-point Gauss--Legendre rule integrates it exactly.

\paragraph{Discontinuous.} By orthogonality, $\int_{[-1,1]^3}P_m^2 =
8/\nu_m$ with $\nu_m=(2a+1)(2b+1)(2c+1)$, so the projection is local:
\begin{equation}
  C_m = \frac{\nu_m}{|V_k|}\int_{V_k}\teps\,P_m .
\end{equation}

\paragraph{Continuous.} The projection solves $\mathsf M\,\mathbf N = \mathbf b$, with
$\mathbf b_I = \int\teps\,\phi_I$ assembled from the cell moments through the
Legendre forms above. On a structured mesh the nodal functions are tensor products,
so $\mathsf M = \mathsf M_x\otimes\mathsf M_y\otimes\mathsf M_z$ with banded 1D mass
matrices, and the solve is three sweeps of 1D solves [7].

\paragraph{Boundary conditions.} Every patch lies in one octant, so it is wholly
incoming or outgoing on each face of the box. For a continuous basis, vacuum is
imposed strongly: the boundary nodal values of incoming patches are zero. This
removes the boundary node from the 1D factor of that axis and patch, and keeps the
Kronecker structure. A reflective boundary would tie a patch to its mirror image
and couple the patches, so it is kept weak: its face residual (below) remains. This
follows the slab's continuous basis (\texttt{bases/continuous\_linear\_space.tex}).

\section{Residual}

Inside cell $k$, energy bin $g$ and patch $j$,
\begin{equation}
  r(\tb,E,\Om) = \sum_m P_m(\tb)\Big[Q_m + T_m(E,\Om) - \Sigt(E)\,C_m
     - \Om\cdot\mathbf G_m\Big],
\end{equation}
where $\mathbf G_m$ are the Legendre coefficients of $\nabla\tpsi$. These follow from
$P_1'=P_0$ and $P_2'=3P_1$ and the factor $2/h$ of each axis. $T_m$ is the in-scatter
of mode $m$. Scattering is local in space, so the transfer moments act mode by mode,
$\Sbar_m = M C_m/(\Delta E_g|A_j|)$, and the tabulated kernel of the correction passes
gives $T_m=\sum C_m U$ with the same table.

In the collision-only iterations, $T_m$ is the binned $\Sbar_m$ and the removal is
the binned $\bar\Sigma_t C_m$. In a correction pass, $T_m$ is the shifted trilinear
interpolant $\hat T_m$ of \texttt{tabulated\_correction.tex}, and $\Sigt(E)$ is
pointwise.

On a face normal to axis $a$, the residual is $-\Omega_a D\,\delta(x_a-x_f)$, with
$D$ the jump of $\tpsi$, a polynomial in the two face coordinates. Its face modes
follow from the traces $P_n(\pm1)=(\pm1)^n$. A continuous basis has no interior
jumps, and none on vacuum faces.

\section{Sampling}

The volume terms (sub-cells of the node grid in $(E,\mu,\varphi)$, one per bin and
patch in the iterations) and the face terms share one CDF over their masses. A point
is uniform in the sub-cell or on the face, with the direction $\propto|\Omega_a|$ on
faces. Its weight is $r/q$, with $r$ evaluated at the point, so the estimate is
unbiased whenever every term with nonzero $r$ has a positive mass. The masses use
\begin{equation}
  \frac{1}{|V_k|}\int_{V_k}|r|\ \le\ \Big(\frac{1}{|V_k|}\int_{V_k}r^2\Big)^{1/2}
  = \Big(\sum_m \frac{c_m^2}{\nu_m}\Big)^{1/2},
\end{equation}
(Cauchy--Schwarz and orthogonality), with $c_m$ the bracket of the residual. This is
evaluated at the sub-cell's $(E,\Om)$ corners, and the largest value is kept, plus
the departure of $\Sigt$ from linear as in \texttt{tabulated\_correction.tex}. The
same root mean square over the face modes gives the face masses. It is positive
wherever the spatial polynomial is not identically zero.

\paragraph{Angular streaming departure (support fix).}
The direction components $\Omega_x=\sqrt{1-\mu^2}\cos\varphi$ and
$\Omega_y=\sqrt{1-\mu^2}\sin\varphi$ are not bilinear. Corner cancellation
therefore does not imply a zero volume residual. On an octant sub-patch each
component lies between its corner extrema, as does its bilinear interpolant.
Add to the corner RMS bound
\begin{equation}
 \sum_{a=x,y}(\max_c\Omega_a-\min_c\Omega_a)
 \left(\sum_m G_{a,m}^2/\nu_m\right)^{1/2}.
\end{equation}
The triangle inequality bounds the spatial RMS departure for every direction;
$\Omega_z=\mu$ has no departure. This is conservative but preserves support,
including in voids. The regression test uses $r=1-\Omega_x-\Omega_y-\Omega_z$
on the positive octant: all corner values vanish while the interior does not.

\section{What is reported}

With an exact residual, the expected update is the projection $\Pi\psi$ of the
true flux; collision-only iterations instead target a binned fixed point. A continuous
projection is global and does not preserve cell averages, so comparing them with
cell tallies of standard Monte Carlo shows a projection error, not a bias
(\texttt{bases/continuous\_linear\_space.tex}). The corrected result therefore adds
the \emph{tallied} cell averages of each correction pass to those of the last
iterate $\tpsi^*$, without projecting them. With an exact residual:
\begin{equation}
  \frac{1}{|V_k|}\int_{V_k}\tpsi^* + \frac{1}{|V_k|}\,\widehat{\int_{V_k}\teps},\qquad
  \E[\cdot\mid\tpsi^*] = \frac{1}{|V_k|}\int_{V_k}\psi .
\end{equation}
This holds because the pass's tally estimates $\int_{V_k}L^{-1}r = \int_{V_k}(\psi -
\tpsi^*)$ without bias. So for every basis the result estimates the true cell averages
without bias when the residual is exact, with the history standard error of the pass.
The implemented scattering correction uses $\hat T$ and retains deterministic
interpolation bias $A L^{-1}(\hat T-T)$ for the cell-average functional $A$.
The history standard error does not include this bias. Node-refinement checks
are required before interpreting it as the total error.

\section{Cost}

Histories per iteration are \texttt{N\_per\_bin} times the number of spatial
unknowns (cell coefficients for DG, nodes for CG), energy bins and patches. Per
cell, DG has 4 or 10 coefficients. CG has about 1 node per cell for $p=1$ and 8 for
$p=2$ on large meshes, more on small ones: the pin's $3\times3\times4$ mesh has 80
and 441 nodes against 36 cells. Higher moments are noisier, since a mode's tally
variance carries the factor $\nu_m$ through the projection.

\section{Verification}

\begin{itemize}
  \item \texttt{test\_space3d.py}:
  \begin{itemize}
    \item each basis reproduces a polynomial of its degree exactly from its cell
    moments, and continuous bases have equal traces across every face;
    \item vacuum boundary nodes of incoming patches are zero after projection;
    \item the segment moments of the tally are exact for all 27 orders.
  \end{itemize}
  \item \texttt{test\_absorber\_box\_polynomial\_bases}: the pure-absorber box of
  \texttt{cartesian\_3d.tex}, with every basis. With no scattering the residual is
  exact, and the corrected cell averages must agree with standard Monte Carlo.
\end{itemize}

\section{Expected limitations}

The residual also contains the variation of the streaming term within a patch,
$(\Om-\bar\Om_j)\cdot\nabla\tpsi$. A higher spatial order does not remove it. In the
slab, with two polar bins, this term was as large as the face residual it replaced
(\texttt{NOTES.md}, section 9). Finer patches, or a basis linear in angle within a
patch, address it.

\begin{thebibliography}{99}
\bibitem{abramowitz} M.~Abramowitz and I.~A.~Stegun, \emph{Handbook of Mathematical
  Functions}, National Bureau of Standards (1964), Ch.~22.
\bibitem{reed} W.~H.~Reed and T.~R.~Hill, ``Triangular mesh methods for the neutron
  transport equation,'' LA-UR-73-479, Los Alamos (1973).
\bibitem{lesaint} P.~Lesaint and P.~A.~Raviart, ``On a finite element method for
  solving the neutron transport equation,'' in \emph{Mathematical Aspects of Finite
  Elements in Partial Differential Equations}, Academic Press (1974).
\bibitem{cockburn} B.~Cockburn, G.~E.~Karniadakis and C.-W.~Shu (eds.),
  \emph{Discontinuous Galerkin Methods: Theory, Computation and Applications},
  Springer (2000).
\bibitem{ciarlet} P.~G.~Ciarlet, \emph{The Finite Element Method for Elliptic
  Problems}, North-Holland (1978).
\bibitem{griesheimer} D.~P.~Griesheimer, W.~R.~Martin and J.~P.~Holloway,
  ``Convergence properties of Monte Carlo functional expansion tallies,''
  \emph{J.\ Comput.\ Phys.} \textbf{211}(1), 129--153 (2006).
\bibitem{vanloan} C.~F.~Van Loan, ``The ubiquitous Kronecker product,'' \emph{J.\
  Comput.\ Appl.\ Math.} \textbf{123}, 85--100 (2000).
\end{thebibliography}

\end{document}
